\uuid{254p}
\exo7id{1580}
\titre{exo7 1580}
\auteur{barraud}
\organisation{exo7}
\datecreate{2003-09-01}
\isIndication{false}
\isCorrection{false}
\chapitre{Réduction d'endomorphisme, polynôme annulateur}
\sousChapitre{Polynôme annulateur}
\cpge{mpsi-mp2i.calcul-matriciel-systemes.e, pcsi.determinants.c, ptsi.determinants.a, mp-mpi.reduction-endomorphismes.g, pc.complements-espaces-vectoriels.d, psi.complements-espaces-vectoriels.d, mp-mpi.reduction-endomorphismes.i, pc.reduction.d, psi.reduction-endomorphismes.d}
\module{Algèbre}
\niveau{L2}
\difficulte{}

\contenu{
\begin{enumerate}
    \item \question{Réduire la matrice 
    $$
    A=
    \begin{pmatrix}
      2&  0& 0\\
      3& -4& 3\\
      3& -6& 5
    \end{pmatrix}
    $$}
    \item \question{Donner un polynôme annulateur de $A$ de degré 2.}
    \item \question{En  déduire qu'il existe des coefficients $a_{n}$ et $b_{n}$ tels que
    $A^{n}=a_{n}A+b_{n}$ et les calculer en fonction de $n$.}
\end{enumerate}
}
